% AUTO-GENERATED by the update_record tool. DO NOT EDIT.
% verify_paper regenerates this file from record.json and the run outputs
% and fails on any difference, so manual edits always surface.
\noindent\textbf{Hypothesis 1.} RIKYU の推論 API（OpenAI 互換）で提供される open-weight モデル 6 種（deepseek-v4.1-flash、glm-5.2、kimi-k2.6、kimi-k3、qwen3.6-35b、qwen3.8-27b）は、SciGym 公開コードの Controller をそのまま用いて SciGym-small 137 件を最大 20 反復で解かせたとき、いずれも平均 Simulation Trajectory Error（STE）が論文 Table 1 の GPT-4.1-mini（0.6007）の水準に届く（0.05 より大きくは上回らない）。~\cite[s1.p1, s1.p5, s1.p10]{duan-2025-measuring}
\begin{enumerate}
\item[\textbf{Claim 1}] deepseek-v4.1-flash で SciGym-small 137 件を最大 20 反復で解かせたときの平均 STE は、論文の GPT-4.1-mini の 0.6007 を 0.05 より大きくは上回らない。~\cite[s1.p2, s1.p3, s1.p4, s1.p5, s1.p10]{duan-2025-measuring}, \cite[s2.p1]{h4duan-2025-scigym}
  \emph{Rationale:} DeepSeek 系で唯一提供されている flash モデルが GPT-4.1-mini の水準に届けば、h1 の deepseek-v4.1-flash の部分を担う。
  \emph{Criterion:} \texttt{\detokenize{deepseek-v4.1-flash}}.\texttt{\detokenize{trajectory_smape}} $-$ 0.6007 $\leq 0.05$.
  \emph{Prediction:} $[-0.3, 0]$ (論文 Table 1 で 2025 年前半の mini モデルは 0.42〜0.63、pro モデルは 0.32〜0.46。2026 年の open-weight thinking モデルは論文の pro と mini の間に入ると見込み、GPT-4.1-mini との差は -0.30〜0 と予測する).
  \emph{Observed:} -0.36024.
  \emph{Verdict:} supported.
\item[\textbf{Claim 2}] glm-5.2 で SciGym-small 137 件を最大 20 反復で解かせたときの平均 STE は、論文の GPT-4.1-mini の 0.6007 を 0.05 より大きくは上回らない。~\cite[s1.p2, s1.p3, s1.p4, s1.p5, s1.p10]{duan-2025-measuring}, \cite[s2.p1]{h4duan-2025-scigym}
  \emph{Rationale:} GLM 系の前世代モデルが GPT-4.1-mini の水準に届けば、h1 の glm-5.2 の部分を担う。
  \emph{Criterion:} \texttt{\detokenize{glm-5.2}}.\texttt{\detokenize{trajectory_smape}} $-$ 0.6007 $\leq 0.05$.
  \emph{Prediction:} $[-0.3, 0]$ (c1 と同じ。前世代のため h1 の中では差が小さい側に出ると見込む).
  \emph{Observed:} pending.
  \emph{Verdict:} pending.
\item[\textbf{Claim 3}] kimi-k2.6 で SciGym-small 137 件を最大 20 反復で解かせたときの平均 STE は、論文の GPT-4.1-mini の 0.6007 を 0.05 より大きくは上回らない。~\cite[s1.p2, s1.p3, s1.p4, s1.p5, s1.p10]{duan-2025-measuring}, \cite[s2.p1]{h4duan-2025-scigym}
  \emph{Rationale:} Kimi 系の前世代モデルが GPT-4.1-mini の水準に届けば、h1 の kimi-k2.6 の部分を担う。
  \emph{Criterion:} \texttt{\detokenize{kimi-k2.6}}.\texttt{\detokenize{trajectory_smape}} $-$ 0.6007 $\leq 0.05$.
  \emph{Prediction:} $[-0.3, 0]$ (c1 と同じ).
  \emph{Observed:} -0.39065.
  \emph{Verdict:} supported.
\item[\textbf{Claim 4}] kimi-k3 で SciGym-small 137 件を最大 20 反復で解かせたときの平均 STE は、論文の GPT-4.1-mini の 0.6007 を 0.05 より大きくは上回らない。~\cite[s1.p2, s1.p3, s1.p4, s1.p5, s1.p10]{duan-2025-measuring}, \cite[s2.p1]{h4duan-2025-scigym}
  \emph{Rationale:} Kimi 系の現行モデルが GPT-4.1-mini の水準に届けば、h1 の kimi-k3 の部分を担う。
  \emph{Criterion:} \texttt{\detokenize{kimi-k3}}.\texttt{\detokenize{trajectory_smape}} $-$ 0.6007 $\leq 0.05$.
  \emph{Prediction:} $[-0.3, 0]$ (c1 と同じ。現行世代の大型モデルのため差が大きい側に出ると見込む).
  \emph{Observed:} -0.375628.
  \emph{Verdict:} supported.
\item[\textbf{Claim 5}] qwen3.6-35b で SciGym-small 137 件を最大 20 反復で解かせたときの平均 STE は、論文の GPT-4.1-mini の 0.6007 を 0.05 より大きくは上回らない。~\cite[s1.p2, s1.p3, s1.p4, s1.p5, s1.p10]{duan-2025-measuring}, \cite[s2.p1]{h4duan-2025-scigym}
  \emph{Rationale:} Qwen 系の 35B モデルが GPT-4.1-mini の水準に届けば、h1 の qwen3.6-35b の部分を担う。
  \emph{Criterion:} \texttt{\detokenize{qwen3.6-35b}}.\texttt{\detokenize{trajectory_smape}} $-$ 0.6007 $\leq 0.05$.
  \emph{Prediction:} $[-0.3, 0]$ (c1 と同じ).
  \emph{Observed:} -0.143611.
  \emph{Verdict:} supported.
\item[\textbf{Claim 6}] qwen3.8-27b で SciGym-small 137 件を最大 20 反復で解かせたときの平均 STE は、論文の GPT-4.1-mini の 0.6007 を 0.05 より大きくは上回らない。~\cite[s1.p2, s1.p3, s1.p4, s1.p5, s1.p10]{duan-2025-measuring}, \cite[s2.p1]{h4duan-2025-scigym}
  \emph{Rationale:} Qwen 系の 27B モデルが GPT-4.1-mini の水準に届けば、h1 の qwen3.8-27b の部分を担う。最も小さいモデルであり h1 の下限を決める。
  \emph{Criterion:} \texttt{\detokenize{qwen3.8-27b}}.\texttt{\detokenize{trajectory_smape}} $-$ 0.6007 $\leq 0.05$.
  \emph{Prediction:} $[-0.3, 0]$ (c1 と同じ。最小のモデルのため差が小さい側に出ると見込む).
  \emph{Observed:} -0.352766.
  \emph{Verdict:} supported.
\end{enumerate}
\noindent\emph{Assumed, so that the claims together imply Hypothesis 1:}
\begin{itemize}
\item temperature 1.0 での 1 回実行の平均 STE（137 件平均）の実行間変動が 0.05 より十分小さいこと。先行する Table 1 再現では件ごとの STE の標準偏差はおよそ 0.27、137 件平均の標準誤差はおよそ 0.02 であった（c1〜c6）
\item RIKYU の推論 API が提供する各モデルが、公開重みと同等の能力を持つこと。量子化や serving の設定は公開されていない（c1〜c6）
\item thinking を出力するモデルでも max\_tokens 32768 の範囲で本文が返り、本文が空のときの呼び直し（最大 3 回）で十分であること（c1〜c6）
\item aarch64 向けの libroadrunner 2.7.0 と antimony 2.14.0 が、論文環境の libroadrunner 2.8.0 と同じ軌道と評価値を与えること（c1〜c6）
\item 平均 STE が GPT-4.1-mini の水準に届くかどうかを代表すること。RMS（modifier あり / なし）と NTS は同じ実行から表として報告するが、判定基準には用いない（c1〜c6）
\item 論文の GPT-4.1-mini の値（0.6007）を比較対象とすること。GPT-4.1-mini を同じ Controller で再実行した先行研究でも 0.60 前後であり、論文値をそのまま使う（c1〜c6）
\end{itemize}
\noindent\textbf{Hypothesis 2.} 論文が報告する「各ファミリーで pro 版が mini 版を一貫して上回る」という構図は、open-weight の GLM-5.3 系でも成り立ち、glm-5.3 は glm-5.3-flash より SciGym-small の平均 STE が低い。~\cite[s1.p6, s1.p10]{duan-2025-measuring}
\begin{enumerate}
\item[\textbf{Claim 7}] glm-5.3 で SciGym-small 137 件を最大 20 反復で解かせたときの平均 STE は、同条件の glm-5.3-flash の平均 STE 以下である。~\cite[s1.p2, s1.p3, s1.p4, s1.p6, s1.p10]{duan-2025-measuring}, \cite[s2.p1]{h4duan-2025-scigym}
  \emph{Rationale:} 同一ファミリー・同世代の大小 2 モデルを同じ Controller で比べる直接の検証であり、h2 の全体を担う。
  \emph{Criterion:} \texttt{\detokenize{glm-5.3}}.\texttt{\detokenize{trajectory_smape}} $-$ \texttt{\detokenize{glm-5.3-flash}}.\texttt{\detokenize{trajectory_smape}} $\leq 0$.
  \emph{Prediction:} $[-0.15, -0.02]$ (論文 Table 1 の pro と mini の STE 差は Gemini で -0.097、GPT で -0.140。open-weight でも同程度か、やや小さい差を予測する).
  \emph{Observed:} pending.
  \emph{Verdict:} pending.
\end{enumerate}
\noindent\emph{Assumed, so that the claims together imply Hypothesis 2:}
\begin{itemize}
\item glm-5.3-flash が glm-5.3 と同世代の小型版であること。両者の構成は公開されておらず、名前から仮定する（c7）
\item 平均 STE の差 1 本で「上回る」を代表すること。論文は STE と RMS の両方で pro > mini としているが、RMS は表で報告するのみで判定には用いない（c7）
\item temperature 1.0 での 1 回実行どうしの比較で、差 0 を境とする判定が実行間変動（標準誤差およそ 0.02）に対して十分頑健であること。差が 0.02 未満のときは結論が揺らぎうる（c7）
\end{itemize}
